% ===========================================================================
%  Apendice A — Fontes Mermaid dos diagramas
% ===========================================================================
\chapter{Fontes Mermaid dos diagramas}
\label{ap:mermaid}

Os diagramas do corpo do documento são vetoriais e desenhados nativamente em
\TeX{}, o que garante tipografia uniforme e nitidez em qualquer ampliação. Este
apêndice fornece o código-fonte \textbf{Mermaid} equivalente de cada um, para uso
direto em \arq{README.md}, plataformas de repositório, \textit{wikis} e
ferramentas de documentação que renderizam Mermaid nativamente.

\begin{nota}
Duas notações do corpo do documento não têm equivalente exato em Mermaid: o
diagrama de casos de uso (Figura~\ref{fig:casosdeuso}) --- Mermaid não implementa
a notação UML de casos de uso, e a aproximação abaixo usa um grafo com
agrupamentos --- e as raias dos diagramas de atividade, expressas aqui como
subgrafos rotulados.
\end{nota}

% ---------------------------------------------------------------------------
\section{Visão geral do sistema}

\begin{lstlisting}[language=mermaid, caption={Visão geral --- fluxo completo entre entrada, aplicação e nuvem}, label={lst:mmvisao}]
flowchart TB
    subgraph ENTRADA["Entrada"]
        SENS["Sensores IIoT / simulador<br/>(replay do dataset)"]
        EV["Novo evento JSON<br/>(REST direto)"]
        USR["Usuario<br/>(chat / novos documentos)"]
    end

    subgraph APP["Estacao de trabalho (docker-compose)"]
        MQTT["Broker MQTT<br/>Eclipse Mosquitto<br/>topic: sensors/+/telemetry"]
        WRK["Ingestion Worker<br/>(paho-mqtt -> valida -> Postgres<br/>-> dispara /diagnose -> alerts)"]
        subgraph UI["Frontend — Streamlit"]
            DASH["Dashboard"]
            CHAT["Chat prescritivo"]
            UPL["Upload de documentos"]
        end
        subgraph API["Backend — FastAPI"]
            RT["/diagnose /chat /events<br/>/stats /documents /health"]
            DIAG["Motor de Diagnostico<br/>KNN + LightGBM"]
            GUARD["Guardrail de cobertura<br/>documental (deterministico)"]
            RAGS["Servico RAG<br/>(retrieve -> prompt -> cite)"]
        end
        PG[("PostgreSQL<br/>historico de eventos")]
        CH[("ChromaDB<br/>chunks dos documentos")]
        ART["Artefatos ML<br/>scaler + knn + lgbm"]
    end

    subgraph GCP["Google Cloud — Vertex AI"]
        EMB["gemini-embedding-2"]
        LLM["gemini-3.6-flash"]
    end

    SENS -->|publish JSON| MQTT
    MQTT -->|subscribe| WRK
    WRK --> PG
    WRK --> RT
    EV --> RT
    USR --> CHAT & UPL
    DASH & CHAT & UPL --> RT
    RT --> DIAG --> GUARD
    GUARD -->|familia documentada| RAGS
    GUARD -->|sem documento| RT
    DIAG <--> PG
    DIAG <--> ART
    RAGS <--> CH
    RAGS --> LLM
    UPL -->|ingestao| CH
    RAGS & UPL --> EMB
\end{lstlisting}

% ---------------------------------------------------------------------------
\section{C4 nível 1 --- Contexto (Figura~\ref{fig:c4l1})}

\begin{lstlisting}[language=mermaid, caption={C4 nível 1 --- contexto do sistema}, label={lst:mmc4l1}]
C4Context
    title Sistema de Manutencao Prescritiva — Contexto

    Person(tec,  "Tecnico de manutencao",   "Diagnostica e executa o procedimento")
    Person(eng,  "Engenheiro de manutencao","Analisa historico e documenta")
    Person(dat,  "Engenheiro de dados",     "Carrega, treina e monitora")

    System(pmx, "Sistema de Manutencao Prescritiva",
        "Identifica o defeito, quantifica ocorrencias semelhantes e prescreve a
         correcao fundamentada na base documental — restrito a falhas documentadas")

    System_Ext(gw,    "Gateway de aquisicao IIoT", "Sensores de vibracao")
    System_Ext(scada, "SCADA / CMMS",              "Ordens de servico")
    System_Ext(vertex,"Vertex AI — Gemini",        "Embeddings, geracao e OCR")

    Rel(tec, pmx, "Diagnostica, consulta")
    Rel(eng, pmx, "Analisa, cadastra documento")
    Rel(dat, pmx, "Carrega, treina")
    Rel(gw,  pmx, "Telemetria", "MQTT")
    Rel(pmx, scada, "Alertas, DLQ", "MQTT")
    BiRel(pmx, vertex, "Embeddings, geracao, OCR", "HTTPS")

    UpdateLayoutConfig($c4ShapeInRow="3", $c4BoundaryInRow="1")
\end{lstlisting}

% ---------------------------------------------------------------------------
\section{C4 nível 2 --- Contêineres (Figuras~\ref{fig:c4l2app} a~\ref{fig:c4l2batch})}

\begin{lstlisting}[language=mermaid, caption={C4 nível 2 --- contêineres (os três planos em uma única vista)}, label={lst:mmc4l2}]
C4Container
    title Sistema de Manutencao Prescritiva — Conteineres

    Person(tec, "Tecnico de manutencao", "")
    Person(eng, "Engenheiro de manutencao", "")
    Person(dat, "Engenheiro de dados / MLOps", "")

    System_Boundary(pmx, "Sistema de Manutencao Prescritiva") {
        Container(ui,   "Interface",           "Streamlit + Plotly", "Paineis, diagnostico, conversa, documentos")
        Container(api,  "API de aplicacao",    "FastAPI + Uvicorn",  "Diagnostico, prescricao, conversa, analitica")
        Container(wk,   "Worker de ingestao",  "paho-mqtt",          "Valida, persiste, diagnostica, alerta")
        ContainerQueue(mq,"Broker MQTT",       "Eclipse Mosquitto 2","telemetry / alerts / dlq")
        Container(scr,  "Scripts de preparacao","Python 3.12 CLI",   "ETL, treino, ingestao documental, simulador")
        ContainerDb(pg, "Historico",           "PostgreSQL 16",      "Eventos, taxonomia, documentos, auditoria")
        ContainerDb(ch, "Base vetorial",       "ChromaDB (HNSW)",    "Chunks + metadados; fonte da cobertura")
        Container(art,  "Artefatos de ML",     "joblib",             "scaler, lgbm, knn, index_meta (74 MB)")
    }

    System_Ext(gw,    "Gateway IIoT", "Publica telemetria")
    System_Ext(scada, "SCADA / CMMS", "Consome alertas")
    System_Ext(vertex,"Vertex AI",    "Gemini")

    Rel(tec, ui, "Usa", "HTTPS")
    Rel(eng, ui, "Usa", "HTTPS")
    Rel(dat, scr,"Executa", "CLI")
    Rel(ui,  api,"Consome", "JSON/HTTP (rede privada)")
    Rel(api, pg, "Le e escreve", "SQLAlchemy")
    Rel(api, ch, "Consulta e indexa", "")
    Rel(api, art,"Carrega em memoria", "joblib")
    Rel(api, vertex, "Embeddings e geracao", "HTTPS")
    Rel(gw,  mq, "Publica telemetria", "MQTT QoS 1")
    Rel(mq,  wk, "Entrega", "MQTT QoS 1")
    Rel(wk,  pg, "merge idempotente", "")
    Rel(wk,  api,"POST /diagnose", "HTTP")
    Rel(wk,  mq, "Publica alerts / dlq", "MQTT QoS 1")
    Rel(mq,  scada, "Entrega alertas", "MQTT")
    Rel(scr, pg, "Carga do historico", "")
    Rel(scr, ch, "Indexacao inicial", "")
    Rel(scr, art,"Gravacao do treino", "")
\end{lstlisting}

% ---------------------------------------------------------------------------
\section{C4 nível 3 --- Componentes da API (Figura~\ref{fig:c4l3})}

\begin{lstlisting}[language=mermaid, caption={C4 nível 3 --- componentes do contêiner de API}, label={lst:mmc4l3}]
C4Component
    title API de aplicacao — Componentes

    Container_Boundary(api, "API de aplicacao (FastAPI)") {
        Component(rt,   "Roteadores REST",        "src/api/main.py",     "/diagnose /chat /stats /events /documents /health /analytics/*")
        Component(dep,  "Injecao de dependencias","src/api/deps.py",     "Sessao, cliente de linguagem, agente — substituiveis em teste")
        Component(mot,  "Motor de diagnostico",   "src/ml/diagnose.py",  "LightGBM + KNN + gate de confianca")
        Component(rag,  "Servico RAG",            "src/rag/service.py",  "Recupera, monta prompt, cita")
        Component(gd,   "Guardrail de cobertura", "src/rag/store.py",    "Decisao deterministica")
        Component(ag,   "Agente e ferramentas",   "src/llm/agent*.py",   "4 ferramentas, memoria de sessao")
        Component(anl,  "Agregacoes analiticas",  "src/api/analytics.py","SQL sob demanda")
        Component(orm,  "Modelos ORM",            "src/core/db.py",      "SQLAlchemy 2.0")
        Component(vs,   "Adaptador vetorial",     "src/rag/store.py",    "ChromaDB")
        Component(cli,  "Cliente de linguagem",   "src/llm/client.py",   "Protocolo substituivel")
        Component(cn,   "Canonizador",            "src/etl/canonize.py", "Regras regex ordenadas")
        Component(bs,   "Bootstrap documental",   "src/rag/bootstrap.py","Indexacao no primeiro boot")
        Component(cfg,  "Configuracao",           "src/core/config.py",  "pydantic-settings")
        Component(sch,  "Esquemas de dominio",    "src/core/schemas.py", "pydantic v2 — UNICOS")
    }

    ContainerDb(pg, "PostgreSQL", "", "")
    ContainerDb(ch, "ChromaDB",   "", "")
    System_Ext(vx,  "Vertex AI",  "")

    Rel(rt, mot, "Diagnostica")
    Rel(rt, rag, "Prescreve")
    Rel(rt, ag,  "Conversa")
    Rel(rt, anl, "Agrega")
    Rel(rt, bs,  "Dispara no startup")
    Rel(dep, cli, "Fornece")
    Rel(mot, rag, "Familia prevista")
    Rel(rag, gd,  "Consulta cobertura")
    Rel(ag,  gd,  "Reutiliza (dentro da ferramenta)")
    Rel(mot, orm, "")
    Rel(rag, vs,  "")
    Rel(gd,  vs,  "")
    Rel(vs,  ch,  "HNSW / cosseno")
    Rel(orm, pg,  "")
    Rel(anl, pg,  "SQL")
    Rel(cli, vx,  "HTTPS")
    Rel(mot, cn,  "")
\end{lstlisting}

% ---------------------------------------------------------------------------
\section{C4 nível 4 --- Classes (Figura~\ref{fig:c4l4})}

\begin{lstlisting}[language=mermaid, caption={C4 nível 4 --- classes do motor de diagnóstico}, label={lst:mmc4l4}]
classDiagram
    class SensorEvent {
        +int id
        +datetime created_at
        +str fault
        +float z_rms_velocity_mm_s
        +float x_rms_velocity_mm_s
        +float temperature_c
        +float ...15 demais grandezas
        +float rpm
        +feature_vector() list~float~
    }
    class DiagnosisEngine {
        -StandardScaler scaler
        -LGBMClassifier clf
        -NearestNeighbors knn
        -DataFrame index_meta
        -LabelCanonizer canonizer
        +diagnose(event) DiagnosisResult
        -_similar_events(family, neighbors)
        -_confidence(probability, agreement) str
    }
    class DiagnosisResult {
        +str family
        +bool is_fault
        +float probability
        +float knn_agreement
        +str confidence
        +SimilarEvents similar
    }
    class SimilarEvents {
        +int count
        +datetime first_seen
        +datetime last_seen
        +float freq_per_day
        +list timeline
        +list neighbor_ids
        +float neighbor_agreement
    }
    class LabelCanonizer {
        -dict families
        -list _rules
        +canonize(raw_label) Family
        +fault_families() list~str~
    }
    class DiagnoseResponse {
        +str predicted_fault
        +bool is_fault
        +float probability
        +float knn_agreement
        +str confidence
        +SimilarEvents similar_events
        +bool documented
        +Prescription prescription
        +str suggestion
    }

    DiagnosisEngine ..> SensorEvent : usa
    DiagnosisEngine ..> DiagnosisResult : produz
    DiagnosisEngine *-- LabelCanonizer
    DiagnosisResult *-- SimilarEvents
    DiagnosisResult ..> DiagnoseResponse : mapeia para
    DiagnoseResponse *-- SimilarEvents
\end{lstlisting}

% ---------------------------------------------------------------------------
\section{Modelo entidade-relacionamento (Figura~\ref{fig:er})}

\begin{lstlisting}[language=mermaid, caption={Modelo entidade-relacionamento}, label={lst:mmer}]
erDiagram
    FAULT_FAMILIES ||--o{ LABEL_MAP    : canonicaliza
    FAULT_FAMILIES ||--o{ EVENTS       : classifica
    FAULT_FAMILIES ||--o{ DOC_COVERAGE : cobre
    FAULT_FAMILIES ||--o{ DIAGNOSES    : preve
    DOCUMENTS      ||--o{ DOC_COVERAGE : referencia
    EVENTS         ||--o{ DIAGNOSES    : gera

    FAULT_FAMILIES {
        int  id       PK
        text name     UK
        bool is_fault
    }
    LABEL_MAP {
        text raw_label PK
        int  family_id FK
    }
    EVENTS {
        bigint      id         PK "identificador de origem, sem autoincrement"
        timestamptz created_at     "indexado"
        text        raw_fault      "identifica tambem a campanha de coleta"
        int         family_id  FK  "indexado, pode ser nulo"
        float       z_rms_velocity_mm_s
        float       x_rms_velocity_mm_s
        float       temperature_c
        float       demais_15_grandezas
        float       rpm
    }
    DOCUMENTS {
        int         id          PK
        text        filename
        text        title
        timestamptz ingested_at
        text        status
    }
    DOC_COVERAGE {
        int family_id   PK
        int document_id PK
    }
    DIAGNOSES {
        bigint      id                   PK
        bigint      event_id                "sem FK: evento ad hoc"
        int         predicted_family_id  FK
        float       probability
        json        neighbors
        json        llm_response
        timestamptz created_at
    }
\end{lstlisting}

% ---------------------------------------------------------------------------
\section{Sequência do diagnóstico (Figura~\ref{fig:seqdiagnose})}

\begin{lstlisting}[language=mermaid, caption={Sequência do diagnóstico completo}, label={lst:mmseqdiag}]
sequenceDiagram
    participant U as Solicitante (tecnico ou worker)
    participant A as API (FastAPI)
    participant D as Motor (LightGBM + KNN)
    participant P as Dados (Postgres + artefatos)
    participant G as Guardrail + RAG (Chroma)
    participant V as Vertex AI (Gemini)

    U->>A: POST /api/v1/diagnose (JSON do evento)
    A->>A: valida esquema SensorEvent
    A->>D: engine.diagnose(event)
    D->>D: 18 brutas + 5 derivadas, padronizadas
    D->>P: predict_proba + kneighbors (k=25)
    P-->>D: familia, probabilidade, 25 vizinhos
    D->>D: concordancia e gate de confianca
    D->>P: estatisticas historicas da familia
    P-->>D: contagem, periodo, freq/dia, serie mensal
    D-->>A: DiagnosisResult

    alt familia de falha COM documento indexado
        A->>G: prescribe(family, event)
        G->>V: embedding + geracao em contexto fechado
        V-->>G: instrucoes em Markdown
        G-->>A: prescricao + citacoes dos metadados
    else familia de falha SEM documento indexado
        A->>G: is_documented(family) -> falso
        Note over G,V: O Vertex AI NAO e chamado.<br/>Aviso + sugestao de registro.
        G-->>A: suggestion
    end

    A-->>U: DiagnoseResponse + auditoria persistida
\end{lstlisting}

% ---------------------------------------------------------------------------
\section{Sequência da ingestão MQTT (Figura~\ref{fig:seqmqtt})}

\begin{lstlisting}[language=mermaid, caption={Sequência da ingestão industrial}, label={lst:mmseqmqtt}]
sequenceDiagram
    participant S as Simulador / gateway
    participant B as Broker (Mosquitto)
    participant W as Worker (paho-mqtt)
    participant P as PostgreSQL
    participant A as API /diagnose

    S->>B: publish sensors/maquina-01/telemetry (QoS 1)
    B->>W: entrega (subscribe sensors/+/telemetry)
    W->>W: valida com o schema pydantic UNICO da API

    alt payload valido
        W->>W: canonicaliza a anotacao, se houver
        W->>P: session.merge(Event) — idempotente por PK
        W->>A: POST /api/v1/diagnose
        A-->>W: diagnostico completo
        W->>B: publish sensors/maquina-01/alerts (se for falha)
        B->>S: entrega ao SCADA / CMMS assinante
    else payload invalido
        W->>B: publish sensors/maquina-01/dlq (motivo + payload)
        Note over W,P: Nada e gravado no banco.<br/>O dado permanece recuperavel na DLQ.
    end
\end{lstlisting}

% ---------------------------------------------------------------------------
\section{Sequência da conversa com agente (Figura~\ref{fig:seqchat})}

\begin{lstlisting}[language=mermaid, caption={Sequência da conversa com o agente}, label={lst:mmseqchat}]
sequenceDiagram
    participant U as Usuario (interface)
    participant E as POST /chat (FastAPI)
    participant G as Agente (LlmAgent)
    participant T as Ferramentas (4 funcoes)
    participant R as Recursos (Postgres, Chroma, Vertex)

    U->>E: pergunta + session_id
    alt agente disponivel
        E->>G: ask(message, session_id)
        G->>G: recupera memoria da sessao (schema adk)
        loop enquanto precisar de ferramenta
            G->>T: decide e invoca ferramenta
            T->>R: SQL ou consulta vetorial
            R-->>T: dados estruturados
            T->>T: registra citacoes no coletor de contexto
            T-->>G: resultado JSON
        end
        G->>R: geracao final em contexto fechado
        R-->>E: texto em Markdown
    else agente indisponivel (sem lib ou credenciais)
        E->>R: RAG de um passo (rag_service.chat)
        R-->>E: texto em Markdown
    end
    E-->>U: resposta + citacoes do coletor + selo de modo
\end{lstlisting}

% ---------------------------------------------------------------------------
\section{Pipeline de dados e documentos (Figuras~\ref{fig:etl} e~\ref{fig:ragpipe})}

\begin{lstlisting}[language=mermaid, caption={Pipeline de dados e pipeline documental}, label={lst:mmpipe}]
flowchart LR
    subgraph ETL["ETL de dados (scripts/ingest_data.py)"]
        CSV["banner.csv<br/>166.796 linhas"] --> CLEAN["Limpeza + dedup +<br/>validacao"]
        CLEAN --> CANON["Canonicalizacao<br/>label_map.yaml<br/>151 rotulos -> 17 familias"]
        CANON -->|rotulo desconhecido| ERR["Excecao: nada e gravado"]
        CANON --> PGL[("PostgreSQL<br/>events + fault_families<br/>+ label_map")]
        PGL --> FEAT["Treino (src/ml/train.py):<br/>StandardScaler + KNN + LightGBM"]
        FEAT --> ARTS["artifacts/<br/>joblib + metrics.json"]
    end
    subgraph DOCPIPE["Ingestao documental (scripts/ingest_docs.py)"]
        PDFS["Doc1..Doc6.pdf<br/>+ novos uploads via API"] --> PARSE["Extracao pymupdf<br/>por secao"]
        PARSE -->|sem camada de texto| OCR["OCR multimodal<br/>(Gemini, 150 dpi)"]
        OCR --> CHUNK
        PARSE --> CHUNK["Chunks ~1200 chars c/ overlap<br/>metadados: doc, secao, pagina, familia"]
        CHUNK --> GEMB["gemini-embedding-2<br/>(dim 768)"]
        GEMB --> CHR[("ChromaDB<br/>collection: manuals")]
    end
\end{lstlisting}

% ---------------------------------------------------------------------------
\section{Casos de uso (Figura~\ref{fig:casosdeuso})}

\begin{lstlisting}[language=mermaid, caption={Casos de uso --- aproximação em grafo (Mermaid não implementa a notação UML de casos de uso)}, label={lst:mmuc}]
flowchart LR
    tec(["Tecnico de manutencao"])
    eng(["Engenheiro de manutencao"])
    dat(["Engenheiro de dados / MLOps"])
    gw(["Gateway de aquisicao IIoT"])
    scada(["SCADA / CMMS"])
    vx(["Vertex AI (Gemini)"])

    subgraph SIS["Sistema de Manutencao Prescritiva"]
        UC01["UC-01 Diagnosticar novo evento"]
        UC02["UC-02 Obter instrucoes de correcao"]
        UC03["UC-03 Tratar falha nao documentada"]
        UC04["UC-04 Registrar documento orientativo"]
        UC05["UC-05 Consultar cobertura documental"]
        UC06["UC-06 Conversar com assistente"]
        UC07["UC-07 Carregar historico de eventos"]
        UC08["UC-08 Treinar motor de diagnostico"]
        UC09["UC-09 Ingerir telemetria industrial"]
        UC10["UC-10 Tratar telemetria invalida"]
        UC11["UC-11 Analisar historico no painel"]
        UC12["UC-12 Verificar saude do sistema"]
    end

    tec --- UC01
    tec --- UC06
    eng --- UC01
    eng --- UC04
    eng --- UC11
    dat --- UC07
    dat --- UC08
    dat --- UC12
    gw  --- UC09
    UC09 --- scada
    UC10 --- scada
    vx  --- UC02
    vx  --- UC04
    vx  --- UC06

    UC01 -.->|include| UC02
    UC03 -.->|extend| UC02
    UC09 -.->|include| UC01
    UC10 -.->|extend| UC09
    UC06 -.->|include| UC05

    classDef sec fill:#f5f5f1,stroke:#c3c2b7;
    class UC03,UC10 sec;
\end{lstlisting}

% ---------------------------------------------------------------------------
\section{Diagramas de atividade (Capítulo~\ref{cap:atividades})}

\begin{lstlisting}[language=mermaid, caption={AT-01 --- carga do histórico corporativo}, label={lst:mmat01}]
flowchart TD
    I(( )) --> A1["Ler o CSV com tipagem de datas"]
    A1 --> A2["Normalizar os rotulos brutos"]
    A2 --> A3["Canonicalizar TODOS os rotulos distintos"]
    A3 --> D1{"Todos casam com alguma regra?"}
    D1 -->|nao| E1["Levantar excecao nomeando o rotulo"]
    E1 --> FE((( )))
    D1 -->|sim| A4["Deduplicar por id e descartar linhas incompletas"]
    A4 --> A5["Criar as tabelas, se ausentes"]
    A5 --> A6["Popular fault_families e label_map"]
    A6 --> A7["Carregar os eventos de forma idempotente"]
    A7 --> A8["Imprimir o resumo da carga"]
    A8 --> F((( )))
\end{lstlisting}

\begin{lstlisting}[language=mermaid, caption={AT-02 --- canonicalização de um rótulo bruto}, label={lst:mmat02}]
flowchart TD
    I(( )) --> A1["Normalizar: minusculas, sem espacos nas bordas"]
    A1 --> A2["Tomar a proxima regra da lista ORDENADA"]
    A2 --> D1{"A expressao casa?"}
    D1 -->|sim| OK["Devolver a familia da regra"]
    OK --> F((( )))
    D1 -->|nao| D2{"Ha mais regras?"}
    D2 -->|sim| A2
    D2 -->|nao| ER["Excecao: rotulo exige curadoria humana"]
    ER --> FE((( )))
\end{lstlisting}

\begin{lstlisting}[language=mermaid, caption={AT-03 --- treino do motor de diagnóstico}, label={lst:mmat03}]
flowchart TD
    I(( )) --> A1["Carregar e canonicalizar o dataset"]
    A1 --> A2["Construir a matriz: 18 brutas + 5 derivadas"]
    A2 --> FORK[" "]
    FORK --> V1["Avaliacao 1<br/>holdout por sessao"]
    FORK --> V2["Avaliacao 2<br/>divisao aleatoria estratificada"]
    FORK --> V3["Avaliacao 3<br/>sessao, sem temperatura"]
    V1 --> JOIN[" "]
    V2 --> JOIN
    V3 --> JOIN
    JOIN --> A3["Retreinar com 100% dos dados"]
    A3 --> A4["Indexar o historico completo no KNN"]
    A4 --> A5["Gravar 4 artefatos + metrics.json"]
    A5 --> F((( )))
\end{lstlisting}

\begin{lstlisting}[language=mermaid, caption={AT-04 --- indexação de um documento orientativo}, label={lst:mmat04}]
flowchart TD
    I(( )) --> A1["Remover chunks anteriores do mesmo arquivo"]
    A1 --> A2["Extrair texto por pagina, rastreando a secao"]
    A2 --> D1{"Algum chunk foi produzido?"}
    D1 -->|nao: escaneado| OCR["Renderizar a 150 dpi e transcrever por OCR"]
    OCR --> A3
    D1 -->|sim| A3["Acumular chunks ~1200 chars, overlap 150"]
    A3 --> A4["Gerar embeddings (RETRIEVAL_DOCUMENT, dim 768)"]
    A4 --> D2{"Contagem de vetores = de textos?"}
    D2 -->|nao| ER["Excecao: desalinhamento silencioso"]
    ER --> FE((( )))
    D2 -->|sim| A5["Indexar um conjunto por familia, com metadados"]
    A5 --> A6["Registrar documento e cobertura no banco"]
    A6 --> F((( )))
\end{lstlisting}

\begin{lstlisting}[language=mermaid, caption={AT-05 --- diagnóstico de um novo evento}, label={lst:mmat05}]
flowchart TD
    I(( )) --> A1["Validar payload contra SensorEvent"]
    A1 --> D0{"Esquema valido?"}
    D0 -->|nao| E0["Responder 422 com os campos inconsistentes"]
    E0 --> F0((( )))
    D0 -->|sim| A2["Calcular as 5 derivadas e padronizar"]
    A2 --> A3["Classificar: familia e probabilidade"]
    A3 --> A4["Recuperar 25 vizinhos e calcular concordancia"]
    A4 --> A5["Agregar estatisticas historicas da familia"]
    A5 --> A6["Classificar a confianca (produto dos dois sinais)"]
    A6 --> D1{"E falha?"}
    D1 -->|nao: estado| A10
    D1 -->|sim| A7["Consultar o guardrail de cobertura"]
    A7 --> D2{"Documentada?"}
    D2 -->|sim| A8["Recuperar, gerar e citar (AT-06)"]
    D2 -->|nao| A9["Aviso + sugestao de registro"]
    A8 --> A11
    A9 --> A11
    A10["Persistir auditoria em diagnoses (best-effort)"] --> A11
    A11["Responder DiagnoseResponse"] --> F((( )))
\end{lstlisting}

\begin{lstlisting}[language=mermaid, caption={AT-06 --- \textit{guardrail} e geração da prescrição}, label={lst:mmat06}]
flowchart TD
    I(( )) --> A1["Obter o nome de exibicao da familia"]
    A1 --> D1{"Existe chunk com esta familia?"}
    D1 -->|nao| N1["Formatar aviso + sugestao de registro"]
    N1 --> FE((( )))
    D1 -->|sim| A2["Montar a consulta: familia + sintomas do evento"]
    A2 --> A3["Gerar embedding (RETRIEVAL_QUERY)"]
    A3 --> A4["Recuperar top-k=6 filtrando pela familia"]
    A4 --> A5["Montar prompt de contexto fechado, trechos numerados"]
    A5 --> A6["Gerar com temperature=0.2 e instrucao restritiva"]
    A6 --> A7["Extrair citacoes dos METADADOS dos trechos"]
    A7 --> F((( )))
\end{lstlisting}

\begin{lstlisting}[language=mermaid, caption={AT-07 --- conversa com o agente}, label={lst:mmat07}]
flowchart TD
    I(( )) --> D0{"Agente disponivel?"}
    D0 -->|nao| FB["Degradar para RAG de um passo"]
    FB --> A5
    D0 -->|sim| A1["Reiniciar o coletor de citacoes"]
    A1 --> A2["Recuperar a memoria da sessao (schema adk)"]
    A2 --> A3["Modelo decide a proxima ferramenta"]
    A3 --> D1{"Qual ferramenta?"}
    D1 --> T1["consultar_historico"]
    D1 --> T2["buscar_documentos (aplica o guardrail)"]
    D1 --> T3["diagnosticar_evento"]
    D1 --> T4["cobertura_documental"]
    T1 --> D2
    T2 --> D2
    T3 --> D2
    T4 --> D2{"Precisa de outra ferramenta?"}
    D2 -->|sim| A3
    D2 -->|nao| A4["Gerar a resposta final em contexto fechado"]
    A4 --> A5["Devolver texto + citacoes + selo do modo"]
    A5 --> F((( )))
\end{lstlisting}

\begin{lstlisting}[language=mermaid, caption={AT-08 --- ingestão de telemetria industrial}, label={lst:mmat08}]
flowchart TD
    I(( )) --> P1["Gateway publica em sensors/{maq}/telemetry (QoS 1)"]
    P1 --> W1["Worker extrai o identificador da maquina do topico"]
    W1 --> D1{"JSON e esquema validos?"}
    D1 -->|nao| DLQ["Compor registro de rejeicao: motivo + payload"]
    DLQ --> B2["Publicar em .../dlq (QoS 1)"]
    B2 --> F1((( )))
    D1 -->|sim| W2["Canonicalizar a anotacao, se houver"]
    W2 --> W3["merge idempotente em events"]
    W3 --> W4["Chamar POST /diagnose"]
    W4 --> D2{"E falha?"}
    D2 -->|sim| B1["Publicar em .../alerts (QoS 1)"]
    B1 --> F2((( )))
    D2 -->|nao| F3((( )))
\end{lstlisting}

\begin{lstlisting}[language=mermaid, caption={AT-09 e AT-10 --- cadastro de documento e indexação automática}, label={lst:mmat0910}]
flowchart TD
    subgraph AT09["AT-09 — cadastro pela interface"]
        I1(( )) --> C1["Usuario abre a secao Documentos"]
        C1 --> C2["Interface lista documentos e cobertura"]
        C2 --> C3["Usuario anexa PDF, titulo e familias"]
        C3 --> CD1{"Campos preenchidos?"}
        CD1 -->|nao| CE1["Erro no formulario"] --> C3
        CD1 -->|sim| C4["POST /documents (multipart)"]
        C4 --> CD2{"Indexacao bem-sucedida?"}
        CD2 -->|nao| CE2["Exibir codigo e detalhe do erro"]
        CD2 -->|sim| C5["Exibir chunks e familias atendidas"]
        C5 --> C6["Invalidar o cache da interface"]
        C6 --> C7["Familia coberta na proxima requisicao"] --> F1((( )))
    end

    subgraph AT10["AT-10 — indexacao automatica no primeiro boot"]
        I2(( )) --> BD0{"RAG_BOOTSTRAP habilitado?"}
        BD0 -->|nao| BN0["Nao fazer nada"] --> F2((( )))
        BD0 -->|sim| BD1{"Indice ja contem familias?"}
        BD1 -->|sim| BN1["Log: bootstrap dispensado"] --> F2
        BD1 -->|nao| BD2{"Indice inspecionavel?"}
        BD2 -->|nao| BN2["Advertir e cancelar"] --> F2
        BD2 -->|sim| B1["Iniciar thread daemon separada"]
        B1 --> B2["Startup conclui: healthcheck passa"]
        B1 --> B3["Indexar cada PDF da cobertura inicial (AT-04)"]
        B3 --> B4["Registrar chunks por documento"] --> F3((( )))
    end
\end{lstlisting}

\begin{lstlisting}[language=mermaid, caption={AT-11, AT-12 e AT-13 --- painel, integração contínua e implantação}, label={lst:mmat111213}]
flowchart TD
    subgraph AT11["AT-11 — consulta a um painel"]
        I1(( )) --> P1["Usuario escolhe secao e subpainel"]
        P1 --> PD0{"API acessivel?"}
        PD0 -->|nao| PE0["Interromper com o endereco configurado"] --> F1((( )))
        PD0 -->|sim| P2["Requisitar APENAS os dados do painel ativo"]
        P2 --> PD1{"Agregacao disponivel?"}
        PD1 -->|nao| PE1["Aviso no painel; demais graficos seguem"] --> P3
        PD1 -->|sim| P3["Aplicar filtros locais"]
        P3 --> P4["Aplicar a paleta validada e o chrome comum"]
        P4 --> P5["Renderizar com legenda explicativa"] --> F2((( )))
    end

    subgraph AT12["AT-12 — integracao continua"]
        I2(( )) --> C1["push em main ou pull request"]
        C1 --> C2["Lint: ruff check + format --check"]
        C2 --> CD1{"Passou?"}
        CD1 -->|nao| CE["Interromper e reportar no PR"] --> F3((( )))
        CD1 -->|sim| C3["Testes: 63, Postgres de servico, LLM fake"]
        C3 --> CD2{"Passou?"}
        CD2 -->|nao| CE
        CD2 -->|sim| C4["Build: docker build da imagem unica"]
        C4 --> CD3{"Passou?"}
        CD3 -->|nao| CE
        CD3 -->|sim| C5["Publicacao autorizada"] --> F4((( )))
    end

    subgraph AT13["AT-13 — provisionamento e implantacao"]
        I3(( )) --> D1["Instalar SDK de IaC e autenticar"]
        D1 --> D2["Inicializar o projeto"]
        D2 --> D3["Planejar: revisar o diff da topologia"]
        D3 --> DD1{"Diff aceitavel?"}
        DD1 -->|nao| DE1["Ajustar a declaracao"] --> D3
        DD1 -->|sim| D4["Aplicar: banco, servicos, volume, grupos"]
        D4 --> D5["Definir segredos por entrada padrao"]
        D5 --> D6["Expor dominios HTTPS e proxy TCP"]
        D6 --> D7["Carregar o historico da maquina local"]
        D7 --> D8["Aguardar a indexacao automatica (AT-10)"]
        D8 --> DD2{"Verificacao pos-implantacao passou?"}
        DD2 -->|nao| DE2["Inspecionar logs e variaveis efetivas"] --> D8
        DD2 -->|sim| D9["Sistema operacional em nuvem"] --> F5((( )))
    end
\end{lstlisting}

\begin{lstlisting}[language=mermaid, caption={AT-14 --- ciclo de tratamento de lacuna documental}, label={lst:mmat14}]
flowchart TD
    I(( )) --> P1["Planta: vibracao anomala detectada"]
    P1 --> S1["Sistema: familia prevista + ocorrencias semelhantes"]
    S1 --> D1{"Documentada?"}
    D1 -->|sim| S2["Prescricao com citacao de doc, secao e pagina"]
    S2 --> T1["Tecnico executa o procedimento"]
    T1 --> T2["Tecnico registra conclusao e valida critérios"]
    T2 --> F1((( )))
    D1 -->|nao| S3["Aviso + sugestao de registro"]
    S3 --> E1["Engenheiro prioriza a lacuna no painel"]
    E1 --> E2["Engenheiro redige o procedimento tecnico"]
    E2 --> E3["Engenheiro cadastra o PDF (AT-09)"]
    E3 --> E4["Familia passa a ser coberta"]
    E4 -->|o proximo diagnostico ja prescreve| S1
\end{lstlisting}

% ---------------------------------------------------------------------------
\section{Infraestrutura (Figuras~\ref{fig:infralocal} e~\ref{fig:infranuvem})}

\begin{lstlisting}[language=mermaid, caption={Implantação local e em nuvem}, label={lst:mminfra}]
flowchart TB
    subgraph LOCAL["Estacao de trabalho — docker compose"]
        LUI["ui :8501<br/>Streamlit"] --> LAPI["api :8000<br/>Uvicorn"]
        LWK["worker<br/>paho-mqtt"] --> LAPI
        LMQ["mosquitto :1883"] --> LWK
        LAPI --> LPG[("postgres :5432")]
        LWK --> LPG
        LAPI --> LVOL[("./data/chroma")]
        LAPI --> LART["./artifacts"]
        LPG --> LPGD[("volume pgdata")]
    end

    subgraph CLOUD["Projeto em nuvem — regiao us-east4"]
        CUI["ui<br/>HTTPS publico"] -->|rede privada| CAPI["api<br/>HTTPS publico"]
        CWK["worker<br/>sem exposicao"] -->|rede privada| CAPI
        CMQ["mosquitto<br/>proxy TCP :1883"] -->|rede privada| CWK
        CAPI -->|rede privada| CPG[("postgres<br/>plugin gerenciado")]
        CWK -->|rede privada| CPG
        CAPI --> CVOL[("volume chroma<br/>1 GB, /app/data/chroma")]
    end

    NET(("Internet")) -->|HTTPS| CUI
    NET -->|HTTPS| CAPI
    NET -->|TCP| CMQ
    CAPI -->|egresso HTTPS| VX["Vertex AI — Gemini"]
    GH["Repositorio Git<br/>push em main"] -.->|build apos CI| CUI
    GH -.->|build apos CI| CAPI
    GH -.->|build apos CI| CWK
\end{lstlisting}

% ---------------------------------------------------------------------------
\section{Agente e ferramentas (Figura~\ref{fig:agente})}

\begin{lstlisting}[language=mermaid, caption={Agente de conversa e suas quatro ferramentas}, label={lst:mmagente}]
flowchart LR
    UI["Streamlit Chat<br/>(session_id por aba)"] --> EP["POST /api/v1/chat"]
    EP -->|ADK disponivel| AG["LlmAgent (gemini-3.6-flash)<br/>Runner + sessoes no Postgres"]
    EP -->|fallback automatico| RAGQ["RAG de um passo<br/>(rag_service.chat)"]
    AG --> T1["consultar_historico<br/>(SQL no Postgres)"]
    AG --> T2["buscar_documentos<br/>(guardrail + ChromaDB)"]
    AG --> T3["diagnosticar_evento<br/>(DiagnosisEngine)"]
    AG --> T4["cobertura_documental"]
    classDef guard fill:#fdede7,stroke:#eb6834;
    class T2 guard;
\end{lstlisting}

% ---------------------------------------------------------------------------
\section{Integração contínua (Figura~\ref{fig:at12})}

\begin{lstlisting}[language=mermaid, caption={Fluxo de integração contínua, forma compacta}, label={lst:mmci}]
flowchart LR
    PUSH["push / PR"] --> LINT["ruff check + format"]
    LINT --> TEST["pytest — 63 testes<br/>(Postgres de servico,<br/>LLM/embeddings mockados)"]
    TEST --> BUILD["docker build"]
    BUILD --> DEPLOY["publicacao autorizada"]
\end{lstlisting}
